\documentclass[12pt,a4paper]{article}
\usepackage{fontawesome5}
\providecommand{\faPencilAlt}{\faIcon{pencil-alt}}
\usepackage{styles/couleurs}
\usepackage{styles/mise_en_page}
\usepackage{styles/boites}
\usepackage{styles/maths}
\usepackage{styles/figures}
\usepackage{enumitem}
\usepackage{array}
\usepackage{tabularx}
\usepackage{amsmath}
\usepackage{amssymb}
\usepackage{tikz}
\usepackage{xparse}
\usetikzlibrary{calc,patterns,arrows.meta}
\setlist[itemize]{leftmargin=1.4em,itemsep=0.15em,topsep=0.2em}
\setlist[enumerate]{leftmargin=*,itemsep=0.25em,topsep=0.25em}
\newcommand{\lignepointillee}{\noindent\dotfill\par\vspace{0.45em}}
\newcommand{\lignesreponse}[1]{%
\begin{tcolorbox}[enhanced,colback=white,colframe=NeutreC!40,boxrule=0.5pt,
arc=1.5mm,left=3mm,right=3mm,top=2mm,bottom=1mm,before skip=2pt,after skip=2pt]
\foreach \i in {1,...,#1}{\lignepointillee}
\end{tcolorbox}}
\newcommand{\titreseq}[1]{%
\begin{center}
{\sffamily\LARGE\bfseries\color{TitreC}Séquence 4 -- Calcul littéral}\\[0.2em]
{\sffamily\large\bfseries\color{SavoirC}#1}
\end{center}}
\begin{document}

\titreseq{Activité 5 -- Réduire une expression}

\remarque{Réduire une expression, c'est l'écrire avec le moins de termes possible
en regroupant les termes de même nature.}

\section*{Partie A -- Des lots identiques}

Un lot contient \(x\) objets.

\begin{enumerate}
  \item Combien d'objets contiennent \(3\) lots ?
  \item Combien d'objets contiennent \(5\) lots ?
  \item Combien d'objets contiennent en tout \(3\) lots et \(5\) lots ?
\end{enumerate}

Complète :

\[
3x+5x=(3+5)x=\ldots
\]

\section*{Partie B -- Trier les termes}

Dans chaque expression, souligne d'une même façon les termes de même nature.

\[
A=3x+7+5x+2
\]

\[
B=4x^2+3x-2x^2+8
\]

\[
C=-5x+4+2x-9
\]

Quels termes peut-on regrouper ? Lesquels faut-il conserver séparément ?

\lignesreponse{5}

\section*{Partie C -- Réduire}

\[
A=3x+5x
\qquad
B=x+x+x+4
\]

\[
C=7x+3-2x+5
\qquad
D=-4x+9+6x-2
\]

\[
E=5x^2+2x-3x^2+7x
\]

\lignesreponse{8}

\section*{Partie D -- Analyser des erreurs}

Un élève écrit \(3x+5=8x\).

Un autre écrit \(2x+3x^2=5x^3\).

Explique les erreurs et corrige.

\lignesreponse{6}

\section*{Partie E -- Vérifier avec une valeur}

Compare les valeurs de \(4x+7-2x+3\) et de \(2x+10\)
pour \(x=5\), puis pour \(x=-3\).

\lignesreponse{5}

\end{document}
